\section{Conclusion}
\label{sec:conclusion}

The first \CodeWM{} paper showed how to keep tiny JEPA code transition
models from collapsing.  This paper asks what the stabilized model is
actually doing internally.  The answer, so far, is that transition signal
lives early: repeated looped encoding compresses it away, explicit
state-space deltas are too small and noisy to serve as the main target,
and predictive deep supervision preserves the intermediate
representations the predictor needs.

The Phase~9 breakthrough is that this is only half of the story.  High
lift can be achieved with a collapsed readout, so the model must also be
checked for rank and cross-example discrimination.  The current practical
recipe is compact: stop the encoder at three loops, add predictive
auxiliary losses on intermediate loops, keep the near-static target, and
use VICReg to keep the pooled readout high-rank while the predictor
learns online-to-target translation.  Against a $125.9$M UnixCoder on a
paper-grade $1000\times5000$ retrieval eval, a $1.3$M \CodeWM{} is
competitive but not uniformly dominant: moderate wins on
\code{by\_joint}, a tie on \code{cos@0.9}, and a slight loss on the
strictest \code{cos@0.95}, at two orders of magnitude less capacity and
an order of magnitude less CPU latency.  Cross-seed runs reveal substantial retrieval variance
($2\text{--}3\times$ KNN@5 gap between peak and next-best seed) and
identify predictor collapse (effective rank $5\text{--}6/128$) as the
next bottleneck, independent of the healthy encoder readout.  The
remaining work is predictor-side regularization, a CodeT5+ retry via
its dedicated embedding head, and a \KernelWM{} transfer test.  The
full-validation
audit is itself a methodological contribution: a batch-128 KNN gallery
inflated our best development-time number by $12\times$, and correcting
it against the $0.02\%$ random baseline surfaced a genuine
$\sim\!20\times$ VICReg-over-baseline improvement in readout
discriminability.  The cross-seed audit (\S\ref{subsec:cross_seed})
adds a second corrective: the headline KNN@5 of $5.80\%$ is a
single-seed peak, not a reproducible point estimate.  Reporting this
honestly, alongside the predictor-collapse finding, sets the next
research direction: the encoder is fixed; the predictor is the
bottleneck.
